\section*{अध्याय \ref{ch:b3:planetary-atmospheres} --- ग्रहीय वायुमंडल और वासयोग्यता}

\begin{solution}{exo:b3:planetary-atmospheres:1}
$H = RT/Mg$: शुक्र $8.314\times737/(0.04401\times8.87) = \qty{15.7}{km}$;
पृथ्वी $8.314\times288/(0.02897\times9.80) = \qty{8.4}{km}$; मंगल
$8.314\times214/(0.04334\times3.71) = \qty{11.1}{km}$; टाइटन
$8.314\times94/(0.0280\times1.35) = \qty{20.7}{km}$। टाइटन का वायुमंडल,
ठंडा होने के बावजूद यह सबसे अधिक विस्तृत है: इसका गुरुत्व
पृथ्वी के गुरुत्व से सात गुना कम है।
\end{solution}

\begin{solution}{exo:b3:planetary-atmospheres:2}
$\lambda = GMm/kTr = \num{6.674e-11}\times\num{5.97e24}\times\num{1.674e-27}
/(\num{1.381e-23}\,T\times\num{6.878e6})$: \qty{800}{K} पर 8.78 तथा
\qty{1600}{K} पर 4.39। पलायन करने वाले अंश हैं $9.78\,e^{-8.78} =
\num{1.5e-3}$ और $5.39\,e^{-4.39} = \num{6.7e-2}$: सौर चक्र
हाइड्रोजन के पलायन अंश को लगभग 44 गुना बदल देता है।
\end{solution}

\begin{solution}{exo:b3:planetary-atmospheres:3}
(क)~$d = S_\mathrm{eff}^{-1/2}$: \omterm{def:b3:planetary-atmospheres:runaway}{आर्द्र ग्रीनहाउस} \qty{0.993}{au}, \omterm{def:b3:planetary-atmospheres:runaway}{अनियंत्रित
ग्रीनहाउस} \qty{0.975}{au}, अधिकतम ग्रीनहाउस \qty{1.705}{au}। (ख)~शुक्र, 1.91 पर,
अनियंत्रित फ़्लक्स से अधिक, और अनुभवजन्य हाल-के-शुक्र सीमा (1.78) से भी
अधिक प्राप्त करता है: यह अंदरूनी किनारे के भीतर पड़ता है। मंगल, 0.43 पर,
1.014 और 0.3438 के बीच पड़ता है, रूढ़िवादी वासयोग्य क्षेत्र के भीतर:
इसकी शुष्कता इसके छोटे आकार और वायुमंडल के नष्ट होने से है,
न कि इसकी दूरी से।
\end{solution}

\begin{solution}{exo:b3:planetary-atmospheres:4}
(क)~$\log_{10}p = 3.9895 - 443.028/93.66 = -0.741$, $p =
\qty{0.18}{bar}$। (ख)~$0.05\times1.6 = \qty{0.080}{bar}$, सापेक्ष
आर्द्रता $0.080/0.18 = 44\,\%$। (ग)~समुद्रों के ऊपर की वायु
संतृप्ति से बहुत दूर है: जैसे पृथ्वी के महासागरों से जल वाष्पित होता है, वैसे ही इनसे मेथेन वाष्पित होती है,
और आर्द्र वायु को बारिश से पहले ऊपर उठना और ठंडा होना पड़ता है।
\end{solution}

\begin{solution}{exo:b3:planetary-atmospheres:5}
$\lambda_\mathrm D = 7.02\times2.014/1.008 = 14.0$। $u_\mathrm D/u_\mathrm H
= \sqrt{1.008/2.014} = 0.71$। प्रति परमाणु अनुपात $= 0.71\times(15.0\,
e^{-14.0})/(8.02\,e^{-7.02}) = 0.71\times\num{1.69e-3} = \num{1.2e-3}$।
\omterm{def:b3:planetary-atmospheres:jeans-escape}{जीन्स पलायन} लगभग कोई ड्यूटेरियम बाहर नहीं जाने देता: यह एक प्रबल छन्ना है, जिसमें
$\alpha \approx 0.001$।
\end{solution}

\begin{solution}{exo:b3:planetary-atmospheres:6}
(क)~$0.016/\num{1.5575e-4} = 103$। (ख)~$\alpha = 0$ होने पर संवर्धन
$M_0/M$ होता है: $M_0 = 103\times\num{6e15} = \qty{6.2e17}{kg}$, या
$\num{6.2e17}/(1000\times\num{4.60e14}) = \qty{1.3}{m}$ जल की परत
शुक्र के ऊपर। (ग)~$M_0/M = 103^{1/0.2} = \num{1.2e10}$, अतीत का भंडार
\qty{7e25}{kg}, पृथ्वी के द्रव्यमान का बारह गुना: असंभव। संवर्धन के लिए
प्रबल प्रभाजन करने वाला पलायन चाहिए (जीन्स या अ-ऊष्मीय),
कम से कम अंतिम जल-हास के दौरान सक्रिय, या
जल की आपूर्ति के साथ स्थायी अवस्था।
\end{solution}

\begin{solution}{exo:b3:planetary-atmospheres:7}
(क)~चार अरब वर्षों में ($\num{1.26e17}$~s) मंगल
$\num{3.8e43}$ परमाणु, $\num{1.9e43}$ जल के अणु, $\qty{5.7e17}{kg}$ खो देता है;
$4\pi(\num{3.3895e6})^2 = \qty{1.44e14}{m^2}$ के क्षेत्र पर यह एक परत है
\qty{3.9}{m} गहरी। (ख)~बीस गुना अधिक, \qty{78}{m}। (ग)~उच्च-तिर्यकता दर भी
केवल 100 से
\qty{1500}{m} के निचले सिरे तक पहुँचती है; शेष भूपर्पटी और हिम में चला गया, या
जब सूर्य युवा था तब तेज़ी से खो गया।
\end{solution}

\begin{solution}{exo:b3:planetary-atmospheres:8}
(क)~$m = pA/g = 636\times\num{1.44e14}/3.71 = \qty{2.5e16}{kg}$।
(ख)~$0.1\times\num{1.26e17} = \qty{1.3e16}{kg}$, वर्तमान वायुमंडल का आधा,
इसका पूरा भाग $\num{2.5e16}/0.1 = \qty{2.5e17}{s}$ में,
\num{7.8}~अरब~वर्ष। (ग)~युवा सूर्य अधिक सक्रिय था और उसके
तूफ़ान अधिक बार आते थे, और तूफ़ानों के दौरान हास बढ़ता है; अन्य मार्ग
(हाइड्रोजन, स्पटरिंग, प्रकाश-रासायनिक पलायन) गिने नहीं गए।
\end{solution}

\begin{solution}{exo:b3:planetary-atmospheres:9}
(क)~$S(1 - 0.193)/4 = 291$ से $S = \qty{1442}{W.m^{-2}}$ प्राप्त होता है, 1.06 गुना
पृथ्वी का। (ख)~आज का शुक्र $0.807\times2601/4 =
\qty{525}{W.m^{-2}}$ अवशोषित करेगा; युवा शुक्र $0.7\times2601 = \qty{1821}{W.m^{-2}}$
प्राप्त करता है, \qty{367}{W.m^{-2}} अवशोषित: दोनों सीमा से ऊपर हैं, इसलिए बादल-रहित युवा शुक्र
शुरू से ही अनियंत्रित हो गया होता। (ग)~बादल
एल्बिडो बढ़ाते हैं; बादलों वाले त्रि-आयामी मॉडलों ने
एक वासयोग्य युवा शुक्र तक का मार्ग पाया है।
\end{solution}

\begin{solution}{exo:b3:planetary-atmospheres:10}
(क)~$H = kT/mg = \num{1.381e-23}\times1000/(16\times\num{1.661e-27}\times9.21)
= \qty{56}{km}$। (ख)~$n_\mathrm{exo} = 1/(\sigma_c H) = 1/(\num{3e-19}\times
\num{5.64e4}) = \qty{5.9e13}{m^{-3}}$; $z = 200 + 56.4\ln(\num{1e16}/
\num{5.9e13}) = 200 + 56.4\times5.13 = \qty{490}{km}$। (ग)~ऊपर गुरुत्व कम होने से
पैमाना-ऊँचाई ऊँचाई के साथ बढ़ती है, इसलिए घनत्व
माने गए से धीरे गिरता है: \omterm{def:b3:planetary-atmospheres:exobase}{बाह्य-आधार} ऊपर खिसक जाता है।
\end{solution}

\begin{solution}{exo:b3:planetary-atmospheres:11}
(क)~$\Phi = \num{3.0e13}\times\num{3e-3} = \qty{9e10}{cm^{-2}.s^{-1}} =
\qty{9e14}{m^{-2}.s^{-1}}$; $4\pi(\num{6.371e6})^2 =
\qty{5.10e14}{m^2}$ के क्षेत्र पर $\num{4.6e29}$ $\ce{H2}$ अणु प्रति सेकंड।
(ख)~एक महासागर $\num{1.4e24}/18.015\times\num{6.022e23} =
\num{4.7e46}$ जल-अणु धारण करता है, प्रत्येक एक $\ce{H2}$ के साथ: यह
$\qty{1.0e17}{s} = \num{3.2}\text{ अरब वर्ष}$ तक टिकता है, जो
पृथ्वी की आयु के तुल्य है, इसी कारण \num{3e-3} का मिश्रण-अनुपात \omterm{def:b3:planetary-atmospheres:runaway}{आर्द्र
ग्रीनहाउस} सीमा को चिह्नित करता है।
\end{solution}

\begin{solution}{exo:b3:planetary-atmospheres:12}
(क)~$p_1 = g/D\kappa = 9.81/(1.66\times0.005) = \qty{1180}{Pa}$।
(ख)~$1/T = 1/273.16 - \ln(1182/611.657)\times461.5/\num{2.5e6}$, $T =
\qty{282.5}{K}$, और $\sigma T^4 = \qty{361}{W.m^{-2}}$। (ग)~वास्तविक सीमा
कम है, इसलिए विकिरण करने वाला स्तर ठंडा है और कम दाब पर है: वास्तविक
जलवाष्प इस धूसर गैस से अधिक अपारदर्शी है, और इसका प्रभावी $\kappa$
बड़ा है।
\end{solution}

\begin{solution}{pb:b3:planetary-atmospheres:1}
\textbf{1.} $g = GM/R^2 = \num{6.674e-11}\times\num{5.2e24}/(\num{6.2e6})^2 =
\qty{9.03}{m.s^{-2}}$; $\bar\rho = M/\tfrac43\pi R^3 = \qty{5210}{kg.m^{-3}}$।

\textbf{2.} $v_\mathrm{esc} = \sqrt{2GM/r} = \sqrt{2\times\num{3.47e14}/\num{6.7e6}}
= \qty{10.18}{km.s^{-1}}$।

\textbf{3.} $u = \sqrt{2kT/m}$: हाइड्रोजन \qty{4.81}{km.s^{-1}}, हीलियम
\qty{2.41}{km.s^{-1}}, ऑक्सीजन \qty{1.21}{km.s^{-1}}।

\textbf{4.} $\lambda = (v_\mathrm{esc}/u)^2$: 4.49, 17.8, 71.2।

\textbf{5.} केवल हाइड्रोजन: उसका पलायित अंश 0.06 है, हीलियम का
\num{3.5e-7}, ऑक्सीजन का \num{9e-30}। इनमें से कोई भी जलगतिक शासन के निकट नहीं है,
$\lambda \approx 1.5$ से 3 तक।

\textbf{6.} $\lambda = 4.49\times1400/3000 = 2.09$: यह संक्रमण में है;
हाइड्रोजन जलगतिक रूप से बहकर निकल गया होता, संभवतः
भारी गैसों को भी अपने साथ खींच ले जाता।

\textbf{7.} $(1 + 4.49)\,e^{-4.49} = 0.0618$।

\textbf{8.} $n u/2\sqrt\pi = \num{2e11}\times4806/3.545 = \qty{2.71e14}{m^{-2}.s^{-1}}$;
$\Phi_\mathrm J = \num{2.71e14}\times0.0618 = \qty{1.68e13}{m^{-2}.s^{-1}}$।

\textbf{9.} $4\pi(\num{6.7e6})^2 = \qty{5.64e14}{m^2}$; $\num{9.5e27}$ परमाणु
प्रति सेकंड।

\textbf{10.} $\num{1.4e24}/18.015 = \qty{7.8e22}{mol}$, $\num{4.68e46}$ अणु,
$\num{9.36e46}$ हाइड्रोजन परमाणु।

\textbf{11.} $\num{9.36e46}/\num{9.5e27} = \qty{9.9e18}{s}$, लगभग 310 अरब
वर्ष।

\textbf{12.} $\lambda_\mathrm D = 4.49\times2.014/1.008 = 8.96$; पलायित अंश
$9.96\,e^{-8.96} = \num{1.28e-3}$; सापेक्ष दक्षता $0.707\times
\num{1.28e-3}/0.0618 = 0.015$।

\textbf{13.} $H_a = kT/m_a g = \num{1.381e-23}\times200/(44.01\times\num{1.661e-27}
\times9.03) = \qty{4.18}{km}$।

\textbf{14.} $b = \qty{1.2e21}{m^{-1}.s^{-1}}$, अतः
\[
\Phi_\mathrm{DL} = \frac{b f(m_a - m_i)}{m_a H_a}
= \frac{\num{1.2e21}\times\num{1e-3}\times(42.0/44.0)}{4185}
= \qty{2.74e14}{m^{-2}.s^{-1}}.
\]

\textbf{15.} $\num{3.0e13}\times\num{1e-3}\times9.03/9.81 =
\qty{2.76e10}{cm^{-2}.s^{-1}} = \qty{2.76e14}{m^{-2}.s^{-1}}$: वही
1\,\% के भीतर।

\textbf{16.} $4\pi(\num{6.2e6})^2 = \qty{4.83e14}{m^2}$; $\num{1.32e29}$ अणु
$\ce{H2}$ के, $\num{2.64e29}$ परमाणु प्रति सेकंड।

\textbf{17.} \omterm{def:b3:planetary-atmospheres:exobase}{बाह्य-आधार} अपने वर्तमान घनत्व पर जितना छोड़ता है, विसरण-सीमा उसका
28 गुना आपूर्ति देती है। तापमंडल में हाइड्रोजन तब तक जमा होता रहता है
जब तक जीन्स फ़्लक्स आपूर्ति के बराबर न हो जाए; दीर्घकाल में पलायन विसरण-सीमित
हो जाता है, जैसा पृथ्वी पर होता है।

\textbf{18.} $\num{2e11}\times28 = \qty{5.6e12}{m^{-3}}$।

\textbf{19.} $\num{4.68e46}/\num{1.32e29} = \qty{3.54e17}{s}$, 11.2 अरब वर्ष।

\textbf{20.} $S = 0.30/0.60^2 = 0.833$।

\textbf{21.} $d = \sqrt{L/S_\mathrm{eff}}$: भीतरी किनारा $\sqrt{0.30/0.939} =
\qty{0.565}{au}$, बाहरी किनारा $\sqrt{0.30/0.290} = \qty{1.02}{au}$। केपोस,
\qty{0.60}{au} पर, भीतर है, भीतरी किनारे के निकट।

\textbf{22.} $0.939/0.833 = 1.13$: 13\,\% की चमक-वृद्धि।

\textbf{23.} $0.833\times1361\times0.7/4 = \qty{198}{W.m^{-2}}$, $T_\mathrm{eq} =
(198/\sigma)^{1/4} = \qty{243}{K}$; \qty{291}{W.m^{-2}} की सीमा से काफ़ी नीचे
(सूर्य-जैसे तारे के लिए परिकलित,
पर अंतर बड़ा है)।

\textbf{24.} तीन गुना तेज़: 3.7 अरब वर्ष।

\textbf{25.} \textbf{11~अरब~वर्ष}।
\end{solution}
